\documentclass[10pt,aspectratio=169]{beamer}
% Preámbulo modular para presentaciones Beamer (Modelación Estadística)
% Diseño y paleta institucional del Tecnológico de Monterrey

\documentclass[aspectratio=169,xcolor={dvipsnames,table}]{beamer}

\usepackage[utf8]{inputenc}
\usepackage[spanish,mexico]{babel}
\usepackage{amsmath,amssymb,amsfonts,mathtools}
\usepackage{graphicx}
\usepackage{listings}
\usepackage{booktabs}
\usepackage{tikz}
\usetikzlibrary{matrix,arrows,decorations.pathmorphing,shapes.geometric,positioning}

% Paleta Institucional Tecnológico de Monterrey
\definecolor{TecAzul}{HTML}{0039A6}       % Azul TEC: color primario institucional
\definecolor{TecAzulOscuro}{HTML}{1A2E51} % Azul marino oscuro
\definecolor{TecRojo}{HTML}{EC2661}       % Rojo de acento (paleta miTec)
\definecolor{TecGris}{HTML}{646464}       % Gris neutro para comentarios y subtítulos
\definecolor{TecFondoBloque}{HTML}{F4F6F9}% Fondo suave para bloques

% Configuración de Tema Beamer
\usetheme{Boadilla}
\usecolortheme[named=TecAzulOscuro]{structure}
\setbeamercolor{alerted text}{fg=TecRojo}
\setbeamercolor{palette primary}{bg=TecAzulOscuro,fg=white}
\setbeamercolor{palette secondary}{bg=TecAzul,fg=white}
\setbeamercolor{palette tertiary}{bg=TecAzulOscuro!80!black,fg=white}
\setbeamercolor{title}{fg=TecAzulOscuro,bg=TecFondoBloque}
\setbeamercolor{frametitle}{fg=TecAzulOscuro,bg=TecFondoBloque}

% Bloques personalizados elegantes
\setbeamercolor{block title}{bg=TecAzulOscuro,fg=white}
\setbeamercolor{block body}{bg=TecFondoBloque,fg=black}
\setbeamercolor{block title alerted}{bg=TecRojo,fg=white}
\setbeamercolor{block body alerted}{bg=TecRojo!8,fg=black}
\setbeamercolor{block title example}{bg=TecAzul,fg=white}
\setbeamercolor{block body example}{bg=TecAzul!8,fg=black}

% Redondear bordes y añadir sombra sutil
\setbeamertemplate{blocks}[rounded][shadow=true]
\setbeamertemplate{navigation symbols}{}

% Pie de página limpio con número de diapositiva
\setbeamertemplate{footline}{
  \leavevmode%
  \hbox{%
  \begin{beamercolorbox}[wd=.7\paperwidth,ht=2.25ex,dp=1ex,left]{author in head/foot}%
    \hspace*{2ex}\insertshortauthor~~(\insertshortinstitute) \hfill \insertshorttitle
  \end{beamercolorbox}%
  \begin{beamercolorbox}[wd=.3\paperwidth,ht=2.25ex,dp=1ex,right]{date in head/foot}%
    \insertshortdate{}\hspace*{2em}
    \insertframenumber{} / \inserttotalframenumber\hspace*{2ex}
  \end{beamercolorbox}}%
  \vskip0pt%
}

% Configuración de Listings de Python para visualización pedagógica de código
\lstset{
    language=Python,
    basicstyle=\ttfamily\footnotesize,
    keywordstyle=\color{TecAzulOscuro}\bfseries,
    stringstyle=\color{TecRojo},
    commentstyle=\color{TecGris}\itshape,
    showstringspaces=false,
    breaklines=true,
    frame=lines,
    rulecolor=\color{TecAzulOscuro!30},
    backgroundcolor=\color{TecFondoBloque},
    aboveskip=0.5em,
    belowskip=0.5em,
    literate={á}{{\'a}}1 {é}{{\'e}}1 {í}{{\'i}}1 {ó}{{\'o}}1 {ú}{{\'u}}1
             {Á}{{\'A}}1 {É}{{\'E}}1 {Í}{{\'I}}1 {Ó}{{\'O}}1 {Ú}{{\'U}}1
             {ñ}{{\~n}}1 {Ñ}{{\~N}}1 {¿}{{?`}}1 {¡}{{!`}}1
}

% Metadatos institucionales por defecto
\title[Modelación Estadística]{Modelación Estadística para Ciencia de Datos e Ingeniería}
\author[J. Castillo Colmenares]{Juliho David Castillo Colmenares}
\institute[Tec de Monterrey]{Tecnológico de Monterrey \\ Escuela de Ingeniería y Ciencias}
\date{\today}

% Comandos y macros estadísticas compartidas para las presentaciones Beamer (Metropolis)

% Conjuntos numéricos básicos
\newcommand{\R}{\mathbb{R}}
\newcommand{\N}{\mathbb{N}}
\newcommand{\Q}{\mathbb{Q}}
\newcommand{\Z}{\mathbb{Z}}
\newcommand{\set}[1]{\left\{ #1 \right\}}
\newcommand{\abs}[1]{\left|#1\right|}
\newcommand{\norm}[1]{\left\| #1 \right\|}

% Comandos estadísticos y probabilísticos del curso
\newcommand{\Var}{\operatorname{Var}}
\newcommand{\cov}{\operatorname{Cov}}
\newcommand{\corr}{\rho}
\newcommand{\comb}[2]{\begin{pmatrix} #1 \\ #2 \end{pmatrix}}
\newcommand{\s}{\sigma}
\newcommand{\particion}{\mathfrak{P}}
\newcommand{\card}[1]{\#\left( #1 \right)}
\newcommand{\seq}[3]{\set{#1}_{#2}^{#3}}
\newcommand{\E}{\mathbb{E}}
\newcommand{\Prob}{\mathbb{P}}

% Entornos pedagógicos y cajas en español académico natural
\newenvironment{discusion}{%
  \begin{alertblock}{Pregunta de Discusión}%
}{%
  \end{alertblock}%
}

\newenvironment{reflexion}{%
  \begin{alertblock}{Reflexión para el Aula}%
}{%
  \end{alertblock}%
}

\newenvironment{paradoja}{%
  \begin{alertblock}{Paradoja de Exploración}%
}{%
  \end{alertblock}%
}

\newenvironment{motivacion}{%
  \begin{exampleblock}{Motivación: Ciencia de Datos en la Práctica}%
}{%
  \end{exampleblock}%
}

\newenvironment{retoclase}{%
  \begin{block}{Reto Rápido en Equipo}%
}{%
  \end{block}%
}

\title[IC y pruebas de hipótesis]{Sección 07.02: IC y pruebas de hipótesis}
\subtitle{De la estimación a la decisión}
\author[J. Castillo Colmenares]{Juliho Castillo Colmenares}
\institute{Tecnológico de Monterrey}
\date{}
\begin{document}
\begin{frame}[plain]\titlepage\end{frame}
\begin{frame}{Hoja de ruta}
\small Esta sección enlaza los intervalos de confianza con las pruebas de dos colas.
\begin{enumerate}\item Regla basada en el valor $p$.\item Equivalencia con un intervalo.\item Lectura e interpretación.\end{enumerate}
\end{frame}
\begin{frame}{Fuente y alcance}
\small La fuente canónica es la sección de teoría sobre la relación entre intervalos de confianza y pruebas de hipótesis. El mazo sigue su orden y conserva sus etiquetas matemáticas.
\begin{block}{Pregunta guía} ¿Cuándo produce la misma decisión un intervalo y una prueba?\end{block}
\end{frame}
\begin{frame}{Decisión estadística}\small
\begin{block}{Regla del valor $p$} Rechazamos $H_0$ cuando $p\leq\alpha$ y no rechazamos $H_0$ cuando $p>\alpha$.\end{block}
\begin{alertblock}{Cuidado} No rechazar $H_0$ no equivale a demostrar que $H_0$ es verdadera.\end{alertblock}
\end{frame}
\begin{frame}{Hipótesis de dos colas}\small
\[H_0:\theta=\theta_0,\qquad H_a:\theta\neq\theta_0.\]
La prueba separa los resultados extremos en dos colas, con una probabilidad total $\alpha$ de rechazo cuando $H_0$ es cierta.
\end{frame}
\begin{frame}{Intervalo de confianza}\small
Un intervalo $(1-\alpha)100\%$ reúne los valores del parámetro compatibles con los datos al nivel elegido.
\begin{block}{Lectura} El intervalo no enumera valores probables del parámetro; describe el procedimiento que captura al parámetro en repeticiones de muestreo.\end{block}
\end{frame}
\begin{frame}{Teorema de equivalencia}\small
\begin{block}{Equivalencia}
Para una prueba bilateral al nivel $\alpha$ y su intervalo de confianza correspondiente,
\[\text{rechazar }H_0\Longleftrightarrow \theta_0\notin IC_{(1-\alpha)100\%}.\]
\end{block}
\end{frame}
\begin{frame}{La misma distancia estandarizada}\small
Ambos procedimientos comparan
\[\frac{\hat\theta-\theta_0}{SE(\hat\theta)}\]
con el valor crítico apropiado, $z_{\alpha/2}$ o $t_{\alpha/2}$.
\end{frame}
\begin{frame}{Caso de una media}\small
Si $\sigma$ es conocida, el intervalo es
\[\bar X\pm z_{\alpha/2}\frac{\sigma}{\sqrt n}.\]
Excluir a $\mu_0$ equivale a que $|\bar X-\mu_0|/(\sigma/\sqrt n)>z_{\alpha/2}$.
\end{frame}
\begin{frame}{Puente computacional I}\small
\begin{itemize}\item Calcular $\hat\theta$ y su error estándar.\item Fijar $\alpha$.\item Obtener el intervalo y el estadístico de prueba.\end{itemize}
\end{frame}
\begin{frame}{Criterio operativo}\small
\begin{enumerate}\item Identificar el parámetro y $\theta_0$.\item Elegir el nivel $\alpha$.\item Construir $IC_{1-\alpha}$.\item Verificar la pertenencia de $\theta_0$.\end{enumerate}
\end{frame}
\begin{frame}{Lectura de una media}\small
\begin{block}{Aplicación derivada de la teoría} Un intervalo para una media tiene centro $\bar x$ y margen $z_{\alpha/2}\sigma/\sqrt n$.\end{block}
\[\mu_0\in IC\Rightarrow\text{no rechazo};\qquad\mu_0\notin IC\Rightarrow\text{rechazo}.\]
\end{frame}
\begin{frame}{Interpretación de una media}\small
La conclusión se expresa respecto de la evidencia y del nivel fijado: el valor hipotético es compatible o no es compatible con el intervalo observado.
\begin{alertblock}{Comunicación} Reportar también extremos del intervalo, unidades y nivel de confianza.\end{alertblock}
\end{frame}
\begin{frame}{Lectura con distribución $t$}\small
Cuando $\sigma$ es desconocida, se sustituye por $S$ y se usa $t_{n-1}$ tanto en el intervalo como en la prueba.
\[\bar X\pm t_{\alpha/2,n-1}\frac{S}{\sqrt n}.\]
\end{frame}
\begin{frame}{Interpretación con $t$}\small
La lógica no cambia: los cuantiles de $t$ reflejan la incertidumbre adicional de estimar $\sigma$. Para muestras grandes, $t_{n-1}$ se aproxima a la normal estándar.
\end{frame}
\begin{frame}{Lectura de dos colas}\small
\begin{block}{Simetría} Una alternativa $H_a:\theta\neq\theta_0$ reparte $\alpha$ entre las dos colas.\end{block}
El intervalo bilateral representa exactamente la región de valores no rechazados.
\end{frame}
\begin{frame}{Aplicación comparativa}\small
Supóngase un intervalo $(1-\alpha)100\%$ para una media. Compare visualmente $\mu_0$ con sus extremos inferior y superior.
\begin{block}{Decisión} La comparación del valor hipotético con los extremos sustituye a dos cálculos de cola.\end{block}
\end{frame}
\begin{frame}{Resolución comparativa}\small
\begin{itemize}\item Si $L\leq\mu_0\leq U$, la prueba bilateral no rechaza.\item Si $\mu_0<L$ o $\mu_0>U$, la prueba rechaza.\end{itemize}
\end{frame}
\begin{frame}{Aplicación de comunicación}\small
Una conclusión útil debe incluir parámetro, intervalo, nivel y decisión. La frase ``no hay evidencia suficiente'' es preferible a ``se acepta la nula''.
\end{frame}
\begin{frame}{Resolución de comunicación}\small
\begin{block}{Plantilla} ``El intervalo de confianza del $100(1-\alpha)\%$ para $\theta$ es $[L,U]$; como $\theta_0$ [pertenece/no pertenece], [no rechazamos/rechazamos] $H_0$.''\end{block}
\end{frame}
\begin{frame}{Síntesis}\small
\begin{itemize}\item La prueba bilateral y el intervalo correspondiente comparten valor crítico.\item La pertenencia de $\theta_0$ determina la decisión.\item El intervalo aporta magnitud, precisión y dirección.\end{itemize}
\end{frame}
\begin{frame}{Conexión con el capítulo}\small
Esta equivalencia permite pasar de estimación a decisión sin duplicar cálculos. Las secciones siguientes aplican el mismo principio a medias, proporciones y varianzas.
\end{frame}
\end{document}
